\documentclass[11pt]{amsart}

\usepackage[margin=0.95in]{geometry}
\usepackage{amsmath,amssymb,amsthm,mathtools}
\usepackage{hyperref}
\usepackage[shortlabels]{enumitem}
\usepackage{xcolor}
\usepackage{unicode-math}

\directlua{
  luaotfload.add_fallback("leanfallback", {
    "NotoSansMath-Regular:mode=harf;range={0x1D400-0x1D7FF}",
    "STIXGeneral-Regular:mode=harf;range={0x1D400-0x1D7FF}",
    "FreeSerif:mode=harf;range={0x1D400-0x1D7FF}"
  })
}
\setmonofont[
  RawFeature={fallback=leanfallback}
]{JuliaMono}

\newtheorem{theorem}{Theorem}[section]
\newtheorem{proposition}[theorem]{Proposition}
\newtheorem{lemma}[theorem]{Lemma}
\newtheorem{corollary}[theorem]{Corollary}
\newtheorem{definition}[theorem]{Definition}
\theoremstyle{remark}
\newtheorem{remark}[theorem]{Remark}
\newcommand{\C}{\mathbb{C}}
\newcommand{\N}{\mathbb{N}}
\newcommand{\abs}[1]{\lvert #1 \rvert}
\newcommand{\norm}[1]{\lVert #1 \rVert}
\DeclareMathOperator{\dist}{dist}
\newcommand{\dd}{\mathrm{d}}
\newcommand{\ee}{\mathrm{e}}
\newcommand{\ii}{\mathrm{i}}
\newcommand{\rhoFn}{\rho}
\newcommand{\Fd}{\mathcal{F}_d}
\newcommand{\Sphere}{S^{2d-1}}
\newcommand{\Ann}[1]{\mathcal{A}_{#1}}

\title{$L^2$-Stability for STFT phase retrieval}
\author{Susanna Bertolini}
\address{Department of Mathematics\\
ETH Z\"urich, Ramistrasse 101, 8092 Z\"urich, Switzerland.
} \email{susanna.bertolini@math.ethz.ch}
\author{Jaume de Dios Pont}
\address{Center for Data Science, New York University, New York, New York 10011, USA}
\email{jaumededios@gmail.com}
\author{Ben Pineau}
\address{Courant Institute for Mathematical Sciences\\
New York University
} \email{brp305@nyu.edu}
\author{Jo\~ao~P.~G.\ Ramos}
\address{Instituto de Matem\'atica Pura e Aplicada (IMPA), Rio de Janeiro, RJ, Brazil}
\email{joao.ramos@impa.br}
\author{Mitchell~A.\ Taylor}
\address{Department of Mathematics\\
ETH Z\"urich, Ramistrasse 101, 8092 Z\"urich, Switzerland.
} \email{mitchell.taylor@math.ethz.ch}

\begin{document}

\begin{abstract}
We prove that the short-time Fourier transform with Gaussian window performs $L^2$-local stable phase retrieval at the constant function. The proof involved significant interplay between mathematicians and LLMs. An autoformalization in Lean 4 of an extension of our result to $L^2$-local stable phase retrieval for all Hermite windows and all elements in the finite span of the canonical basis vectors is also presented. 
\end{abstract}

\maketitle

\section{Introduction}
Short-time Fourier transform (STFT) phase retrieval refers to the problem of recovering an unknown function $f\in L^2(\mathbb{R}^d)$ from the phaseless measurement $|V_gf|$, where $V_g:L^2(\mathbb{R}^d)\to L^2(\mathbb{R}^{2d})$ denotes the STFT with window function $g\in L^2(\mathbb{R}^d)$. This problem arises in ptychography, which is a widely used imaging technique that obtains information about an unknown object $f$ by illuminating it with shifts of a probe $g$. 
Due to the high oscillation of the diffracted waves, practitioners can only record the far-field intensity $|V_gf|^2$, which leads to the phase retrieval problem described above.
\medskip

The mathematics of STFT phase retrieval is a rich subject, with contributions coming from many different communities. Recall that, given a window function $g\in L^2(\mathbb{R}^d)$, the \emph{short-time Fourier transform} $V_g:L^2(\mathbb{R}^d)\to L^2(\mathbb{R}^{2d})$ is defined by
\begin{equation*}\label{STFT Def}
V_g f(x,\omega)=\mathcal{F}(T_x\overline{g}f) (\omega)=\int_{\mathbb{R}^d} f(t)\overline{g(t-x)}e^{-2\pi it\cdot \omega}dt, \ \ \text{for}\ x,\omega\in \mathbb{R}^d.
\end{equation*}
It is easy to see that 
\begin{equation}\label{isometry}
    \|V_g f\|_{L^2(\mathbb{R}^{2d})}=\|g\|_{L^2(\mathbb{R}^{d})} \|f\|_{L^2(\mathbb{R}^{d})}
\end{equation} and the fundamental problem in STFT phase retrieval is to stably reconstruct $f\in L^2(\mathbb{R}^d)$ from the phaseless measurement $|V_g f|$, up to global phase ambiguity. \medskip

If stability issues are ignored, there is a standard method for determining whether the operator $V_g$ does phase retrieval. Recall that the \emph{ambiguity function} $\mathcal{A}f$  of a square-integrable function $f$ is defined by $\mathcal{A}f(x,\omega):=e^{\pi ix\cdot \omega} V_ff(x,\omega)$. It is well-known that $\mathcal{A}f$ is a continuous function that determines every $f\in L^2(\mathbb{R})$ up to global phase. Moreover, we have the relation 
\begin{equation*}\label{Ambiguity to STFT}
    \mathcal{F}\left( |V_g f|^2\right)(\omega,-x)=\mathcal{A}f(x,\omega) \overline{\mathcal{A}g(x,\omega)}.
\end{equation*}
For the Gaussian window $g$, it is easy to see that $\mathcal{A}g$ is again a Gaussian, so the above identity implies that $V_g$ does phase retrieval. Similarly, since the ambiguity function of each Hermite window only vanishes on circles, these windows  do STFT phase retrieval
\medskip

The stability problem for STFT phase retrieval is fundamentally difficult and has been studied by many authors. The most natural notion of local stability is as follows.
\begin{definition}\label{defn of local stability}
    Fix a window function $g\in L^2(\mathbb{R}^d)$. We say that $V_g$ does \emph{local stable phase retrieval at $f\in L^2(\mathbb{R}^d)$} if there exists a constant $C(f)<\infty$ such that for all $h\in L^2(\mathbb{R}^d)$ we have
    \begin{equation}\label{SPR local V_g}
        \inf_{|\lambda|=1}\|f-\lambda h\|_{L^2(\mathbb{R}^d)}\leq C(f) \| |V_g f|-|V_g h|\|_{L^2(\mathbb{R}^{2d})}.
    \end{equation}
\end{definition}
Despite extensive efforts on this problem, we are unaware of a single pair of functions $f,g\in L^2(\mathbb{R}^d)\setminus\{0\}$ for which $V_g$ does local stable phase retrieval at $f$. On the other hand, for \emph{any} window $g\in L^2(\mathbb{R}^d)\setminus\{0\}$, it is known \cite{alaifari2025cheeger,Alharbi2024Locality} that one can construct a dense set of functions $f\in L^2(\mathbb{R}^d)$ for which $C(f)=\infty$. 
\medskip

The objective of the present article is to show that there are many functions $f,g\in L^2(\mathbb{R}^d)\setminus\{0\}$ for which $V_g$ does local stable phase retrieval at $f$ in the sense of Definition~\ref{defn of local stability}. In fact, when $g$ is the Gaussian or any other Hermite window, there is a  canonical dense set of functions with $C(f)<\infty.$


\subsection{History and preliminaries} Influential results of Cahill-Casazza-Daubechies \cite{MR3554699} and Alaifari and Grohs \cite{MR3656501} show that phase retrieval using frames  cannot be uniformly stable, unless the Hilbert space $\mathcal{H}$ is finite-dimensional. More precisely,  whenever $T:\mathcal{H} \to L^2(\mu)$ is the analysis operator of a frame and $\dim\mathcal{H}=\infty$, we have $\sup_{f\in \mathcal{H}}C(f)=\infty,$ where $C(f)\in [0,\infty]$ is the smallest constant satisfying
\begin{equation}\label{LSPR intro}
    \inf_{|\lambda|=1}\|f-\lambda g\|_{\mathcal{H}}\leq C(f)\| |Tf|-|Tg|\|_{L^2(\mu)} \ \text{for all} \ g\in \mathcal{H}.
\end{equation} 
 In fact, it was recently shown in \cite{Alharbi2024Locality} that we have $C(f)=\infty$ for a dense collection of $f\in \mathcal{H}$. In particular, in the natural setting of Definition~\ref{defn of local stability}, instabilities cannot be avoided. 
 \medskip

 In view of the above negative results, the question of \emph{how} instabilities occur became fundamental. Importantly, the instabilities found in the above works can be viewed as mathematical artifacts rather than inherent issues with the experimental setup. Indeed, they all occur (in some sense) ``at infinity". This intuition has been made precise in the celebrated works of Alaifari-Daubechies-Grohs-Yin \cite{alaifari2019Stable} and Grohs-Rathmair \cite{Grohs2019Stable,MR4404785}, and more recently in \cite{alaifari2025cheeger}. However, all of these papers have the significant drawback that they only apply to the Gaussian window and (for technical reasons) they must replace the natural $L^2$-norm on the right-hand side of \eqref{SPR local V_g} with a weighted first order Sobolev norm.
 \medskip
 
 The lack of robustness when the window is changed is obviously an important issue, as the window is not precisely known in practice. However, the issue of changing the topology is equally problematic, as it causes the forward problem $f\mapsto |V_gf|$ to be ill-posed. In view of this, it has been an important open problem to understand when stability holds in the natural $L^2$-norm and to what extent one may replace the Gaussian window by a more general window \cite{rathmair2024stable}.

\subsubsection{Reformulating the problem}
Let $g\in L^2(\mathbb{R}^d)$ satisfy $\|g\|_{ L^2(\mathbb{R}^d)}=1$. In view of the identity \eqref{isometry}, Definition~\ref{defn of local stability} can be seen as a special case of the following unified problem by taking $E=V_g(L^2(\mathbb{R}^{d}))\subseteq L^2(\mathbb{R}^{2d})$.

\begin{definition}\label{defn of local stability2}
    Let $E$ be a closed subspace of $L^2(\Omega,\mu)$ for some measure space $(\Omega,\mu)$. We say that $E$ does \emph{local stable phase retrieval at $f\in E$} if there exists a constant $C(f)<\infty$ such that for all $h\in E$ we have
    \begin{equation*}
        \inf_{|\lambda|=1}\|f-\lambda h\|_{L^2(\mu)}\leq C(f) \| |f|-|h|\|_{L^2(\mu)}.
    \end{equation*}
\end{definition}
In a pioneering work by Calderbank-Daubechies-Freeman-Freeman \cite{calderbank2022stable}, infinite-dimensional subspaces $E$ of real-valued $L^2(\mathbb{R}^d)$ have been constructed that satisfy $\sup_{f\in E}C(f)<\infty$. Similar examples over the complex field were later constructed in \cite{MR4837558} and a general theory of stable phase retrieval in Banach lattices was initiated in \cite{MR4800909}. This has led to a host of new and exciting results \cite{camunez2025characterization,garcia2025isometric,garcia2025existence} that lie in the intersection of various fields. When we state our main results below, we will use the language of Definition~\ref{defn of local stability2} rather than Definition~\ref{defn of local stability}.

\subsubsection{Function spaces} 

For any measurable function $F \colon \mathbb{C} \to \mathbb{C}$, we define the
\emph{Fock norm}
\[
\lVert F \rVert_{\mathcal{F}}^2
\;:=\lVert F \rVert_{\gamma}^2
\;:=\;
\frac{1}{\pi}\int_{\mathbb{C}} \lvert F(z) \rvert^2 \, \mathrm{e}^{-\lvert z \rvert^2}\,\mathrm{d} m(z).
\]
The \emph{Bargmann--Fock space}~$\mathcal{F}^2(\mathbb{C})$ is the closed subspace of entire functions with
$\lVert F \rVert_\mathcal{F} < \infty$.
The monomials $e_n(z) := z^n/\sqrt{n!}$ form an orthonormal basis of $\mathcal{F}^2(\mathbb{C})$, so every
$F(z) = \sum_{n \ge 0} a_n z^n$ satisfies
$\lVert F \rVert_{\mathcal{F}}^2 = \sum_{n \ge 0} \lvert a_n \rvert^2 n!$.
With this convention, the Fock norm is defined for arbitrary measurable
functions --  not only holomorphic ones -- and in particular we have
$\lVert \,\lvert F \rvert\, \rVert_\mathcal{F} = \lVert F \rVert_\mathcal{F}$.
\medskip

Up to universal factors, one may identify the Fock space with the image of the short-time Fourier transform with Gaussian window.
Since elements of the Fock space are entire functions, it is obvious that phase retrieval holds in this subspace, i.e., for all $F,G\in \mathcal{F}^2(\mathbb{C})$ we have $|F|=|G|$ if and only if $F=\lambda G$ for some unimodular scalar $\lambda$. However, the question of stability is far more subtle. In fact, it was recently shown that there is a dense set of elements of $\mathcal{F}^2(\mathbb{C})$ that \emph{fail} local stable phase retrieval (even with derivative loss). The main goal of this article is to construct a complementary dense set that \emph{does} local stable phase retrieval in the natural $\mathcal{F}$-norm.
\medskip

Our stability results  apply not only to the Gaussian window, but to any Hermite window. In this case (again up to universal factors), the STFT maps $L^2(\mathbb{R})$ onto the \emph{true polyanalytic Fock space} $\mathcal{H}_k(\mathbb{C})$. Here, for any integer $k\geq 0$, we define $\mathcal{H}_k(\mathbb{C})$ as the $\|\cdot\|_\gamma$-closed linear span of the true Hermite basis
\begin{equation*}
    \Phi_{k,n}(z)=\frac{1}{\sqrt{k!n!}}\sum_{j=0}^{\min(k,n)}(-1)^j{k\choose j}\frac{n!}{(n-j)!}z^{n-j}\overline{z}^{k-j}, \ \ n\geq 0.
\end{equation*}
Clearly, $\mathcal{H}_0(\mathbb{C})=\mathcal{F}^2(\mathbb{C})$ is the Bargmann--Fock space and $\mathcal{H}_k(\mathbb{C})$ is a subspace of polyanalytic functions. By \eqref{Ambiguity to STFT}, $\mathcal{H}_k(\mathbb{C})$  performs phase retrieval. However, the major conceptual difference compared with the case $k=0$ is that the ambiguity functions of the higher Hermite windows have non-trivial zero sets. A priori, this could effect the stability of the phase recovery process, since solving for $\mathcal{A}f$ in \eqref{Ambiguity to STFT} would require division by the ambiguity function of the window.
\subsection{Main results}
The main results of this paper concern the $L^2$-local stability of STFT phase retrieval for the canonical window functions at elements of the span of the canonical basis vectors. However, rather than stating one main result, we prefer to state the result in stages, following the lean verification below.
\medskip

The main result for which we will present a detailed proof is the following.
\begin{theorem}[Local SPR at $\mathbf{1}$ in $\mathcal{F}^2(\mathbb{C})$]\label{main:local}
There exists an absolute constant $M> 0$ such
that for any $F\in \mathcal{F}^2(\mathbb{C})$ we have
\[
\inf_{|\lambda|=1}\lVert F - \lambda \rVert_\mathcal{F}
\;\le\;
M\,
\bigl\|\, \lvert F \rvert - 1 \,\bigr\|_{\mathcal{F}}.
\]
\end{theorem}


Theorem~\ref{main:local} follows from Theorem~\ref{thm:local}, a standard density argument, and a general fact about reproducing kernel Hilbert spaces that we will discuss below. 
\begin{theorem}[Local-local SPR at $\mathbf{1}$ in $\mathcal{F}^2(\mathbb{C})$]\label{thm:local}
There exist absolute constants $\delta > 0$ and $\widetilde{M} > 0$ such
that the following holds.  Let $p$ be a polynomial with
$\lVert p \rVert_\mathcal{F} \le \delta$.  Then there exists $w \in \mathbb{C}$ with $\lvert w \rvert = 1$
such that
\[
\lVert w(1+p) - 1 \rVert_\mathcal{F}
\;\le\;
\widetilde{M}\,
\bigl\|\, \lvert 1+p \rvert - 1 \,\bigr\|_{\mathcal{F}}.
\]
One may take $\delta = 1/4601$ and $\widetilde{M} = 23003$.
\end{theorem}
\begin{remark}
    Theorem~\ref{thm:local} has been autoformalized in Lean 4.  We refer the reader to Subsection~\ref{LLM Sect} for a discussion of the lean formalization and to \cite{armstrong2026formalization} for recent efforts towards autoformalization in analysis. We remark that the assumption that $p$ is a polynomial in Theorem~\ref{thm:local} is for convenience, as a standard density argument allows one to pass from a polynomial $p$ to a general element $F\in \mathcal{F}^2(\mathbb{C})$. The reduction from Theorem~\ref{thm:local} to Theorem~\ref{main:local} is a general RKHS fact that does not rely on any further structure of the Fock space, so will be Lean verified later in a more systematic way -- we include the complete natural language proof in Section~\ref{subsec:local-reduction}.
\end{remark}

Theorem~\ref{main:local} can be vastly generalized using similar methods. Since the proof becomes much more technical but not more instructive, we will only sketch the main changes and provide a link to the full Lean verification (again in the style of Theorem~\ref{thm:local}) together with the natural language proof used to generate the Lean code. The most general result we have in dimension one establishes local stability for every Hermite window at every element in the linear span of the canonical basis vectors.
\begin{theorem}[General stability for Hermite phase retrieval]\label{thm:local_Gen}
Let $k\geq 0$ and $F_0\in \text{span}\{\Phi_{k,n}\}_{n=0}^\infty$. There exists a constant $M=M(k,F_0)>0$ such
that for any $G\in \mathcal{H}_k(\mathbb{C})$ we have
\[
\inf_{|\lambda|=1}\lVert F_0 - \lambda G \rVert_\gamma
\;\le\;
M\,
\bigl\|\, \lvert F_0 \rvert - |G| \,\bigr\|_{\gamma}.
\]
\end{theorem}

\begin{remark}
    The above results have higher dimensional variants; we refer the reader to the associated GitHub for details.
\end{remark}

As a direct consequence of Theorem \ref{main:local}, we are able to prove a \emph{sharp} stability version of the log-Sobolev inequality for entire functions. More specifically, if $F \in \mathcal{F}$ with $\|F\|_{\mathcal{F}}= 1$ the log-Sobolev inequality for entire functions states that 
\[
\|F'\|_{\mathcal{F}} \geq \||F|^2 \log |F|\|_{\gamma}.
\]
This estimate is directly tied to the classical logarithmic Sobolev inequality (see e.g. \cite{Gross1975,Stam1959,Weissler1978} for the first works dealing with such an estimate), and its \emph{stability} has been a recent object of interest of several authors. 

As a matter of fact, we first highlight \cite{DolbeaultEstebanFigalliFrankLoss2025}, where the authors prove asymptotically sharp stability estimates for the Sobolev inequality, which in turn yield, by a limiting argument, stability estimates for the \emph{classical} log-Sobolev inequality. More recently, however, R. Frank, F. Nicola and P. Tilli \cite{FrankNicolaTilli2025} showed that the same estimate holds also in the aforementioned \emph{entire} case, obtaining the following result: 

\begin{theorem}[Frank--Nicola--Tilli, {\cite[Thm.~14]{FrankNicolaTilli2025}}]\label{thm:FNTLogSob}
There exists a constant $c_\ast > 0$ such that, for every
$F \in \mathcal{F}$ with $\|F\|_{\mathcal{F}} = 1$,
\begin{equation}
  \int_{\mathbb{C}} |F'(z)|^2 \, d\gamma(z)
  \;\ge\;
  \int_{\mathbb{C}} |F(z)|^2 \ln |F(z)|^2 \, d\gamma(z)
  \;+\;
  c_\ast \!\!\inf_{\substack{\beta \in \mathbb{C} \\ |c|=1}}\!
  \bigl\| F - c\, e^{\beta z - |\beta|^2/2} \bigr\|_{\mathcal{F}}^{2}.
\end{equation}
\end{theorem}

As we shall see, Theorem \ref{thm:FNTLogSob} follows at once from Theorem \ref{main:local} and \cite[Corollary~1.2]{DolbeaultEstebanFigalliFrankLoss2025}, highlighting how powerful Fock-space phase retrieval is. 

\subsection{Explanation of LLM usage and Lean}\label{LLM Sect}
Here we explain the events that led to the results in this paper.
\medskip

\textbf{LLM Usage.} The problem under consideration had previously been considered by several of the authors, with varying degrees of success. On 28 March 2026, the authors were experimenting with ChatGPT Plus and became intrigued by a (false) reduction of the problem to a circle estimate (which one may view as a primitive version of the results in Section~\ref{subsec:circle}). The authors then tried to understand this reduction, further prompting GPT until we realized that a modification of the proposed argument could be salvaged and would likely lead to a proof. We then sketched the full proof of Theorem~\ref{thm:orthog} by ourselves, which naively lost a $\log$-factor in the stability estimate, and asked Codex to refine it. This ultimately resulted in a complete proof of Theorem~\ref{thm:orthog}, which we knew implied Theorem~\ref{main:local}. We then autoformalized the result in Lean 4 using Claude Opus 4.6.
\medskip


In hindsight, the fact that the structure of the problem is more easily visible in polar coordinates should have been evident. In fact, spiral sampling is commonly used in experimental applications. Moreover, the final proof is not difficult, and once the structure was made clear, the authors could have easily completed the proof of  Theorem~\ref{main:local}  without any additional assistance. What is interesting, however, is that the LLMs were able to incrementally generalize Theorem~\ref{main:local} to higher dimensional variants of Theorem~\ref{thm:local_Gen} with very little human effort. Of course, it should be emphasized that the authors strongly believed that the method would generalize. Moreover, as we will see below, significant challenges arose on the Lean side when verifying the full result.
\medskip

As alluded to above, several of the intermediate results in this paper were known to the authors for a long time. For example, the ``local-local to local" reduction in Theorem~\ref{prop:local-reduction} is relatively standard. The orthogonal reduction in Theorem~\ref{thm:reduction}  is also just a natural variant of \cite[Theorem 3.10]{MR4800909} and was not produced by an LLM.
\medskip

\textbf{Lean.} The main formal content of the results in this paper have been verified in Lean 4. Here, we point out some interesting features of the formalization that may be useful for other problems. We refer the reader to Section~\ref{App} for a formalized statement of our main theorem.
\medskip

First, we carried out the Lean verification only for finite linear combinations of the basis elements. Our justification for this was threefold. First, the density argument is completely standard -- see Section~\ref{subsec:local-reduction} for details. Second, the density argument is a general fact that we plan to formalize in a systematic way. Finally, we found that working with arbitrary polyanalytic functions rather than ``finite" polynomials posed significant difficulties in the Lean formalization. 
\medskip


The formalization process was completed with extensive assistance from LLMs; namely, GPT 5.4 and Claude Opus 4.6, under the supervision of the authors. The authors checked the statements of the main theorems and provided support whenever the formal development encountered blockers, such as statement mismatches between the blueprint and the corresponding results in the Lean files, or overly strong statements that made the proof substantially harder to fill in. We found, for example, that it was worthwhile to make all statements as quantitative as possible and to avoid all compactness arguments -- remnants of this insistence can be (intentionally) witnessed by the explicit choice of constants in the proof, which we did not try to optimize. The full Lean code is available at \textcolor{red}{???}. 
\medskip

The main difficulties of the formalization arose from incorrect translations from the blueprint to the scaffolding. This led to incorrect or imprecise scaffolding files even with an extremely detailed blueprint. Moreover, another recurring issue was the tendency of the LLMs to incorporate unnecessary Hilbert-space structure whenever the blueprint mentioned it for context, even if the main theorem was clearly stated for finite linear combinations only. This introduced many density arguments that did not play a role in the proof of the main theorem but created blockers that had to be recognized as irrelevant and be removed before being able to complete the formalization.



\subsection{Acknowledgments}
The authors thank Joan Bruna, Javier Gomez-Serrano, Alessio Figalli, and Joaquim Serra for their support and guidance during the preparation of this article. J.dD.P and M.A.T.~also thank Manuel Guizar Sicairos and Abraham Levitan for hosting them at PSI and for the instructive discussions on the experimental aspects of phase retrieval. Finally, we thank the Banach lattice and phase retrieval communities for their encouragement to pursue this line of research.

\section{Proof of Theorem~\ref{main:local}}
Here, we present the proof of Theorem~\ref{main:local} and describe the main technical changes which permit a generalization to local stable phase retrieval for all Hermite windows at all functions in the finite span of the canonical basis vectors. 

\subsection{Outline of the proof}\label{sec:strategy}
Below, we give an outline of the logical structure of the proof of Theorem \ref{main:local} as well as a brief overview of the key ideas. We will elaborate on each point in more detail in the corresponding subsection. 
\medskip

\noindent\textbf{The local reduction
(Section~\ref{subsec:local-reduction}).}\;
The main goal of this subsection will be to show that Theorem \ref{main:local} is implied by Theorem \ref{thm:local}. This will follow from a more general argument that holds in any reproducing kernel Hilbert space (which the Fock space is a particular example of). Essentially, we will show that if one
can prove phase retrieval at a function \(F_0\) and establish local stability
with respect to very small perturbations of \(F_0\), then one can prove local
stability for general perturbations of \(F_0\). In the context of Theorem \ref{main:local}, one takes $F_0=1$, but we prefer to state this argument in a more abstract form so that it also applies (for instance) in the context of Theorem \ref{thm:local_Gen} and its higher dimensional analogues. 
\medskip

\noindent
\textbf{The orthogonal reduction (Section~\ref{subsec:reduction}).}\;
Another general argument will be used to show that if orthogonal perturbations
of a unit vector $F_0$ satisfy a certain coercivity estimate with constant~$M$,
then all sufficiently small perturbations (after optimizing over a global
phase) satisfy the same coercivity estimate with a comparable constant $\widetilde{M} = O(M)$. We leave the precise statement and proof of this reduction to Section \ref{subsec:reduction}, but we emphasize that it is just the natural local variant of \cite[Theorem 3.10]{MR4800909}. In the special case of the Fock space $\mathcal{F}^2(\mathbb{C})$, this observation will enable us to reduce the proof of Theorem
\ref{thm:local} to the following ``orthogonal
coercivity" estimate, whose proof is the technical heart of the paper.

\begin{theorem}[Fock-space coercivity]\label{thm:orthog}
There exists an absolute constant $C$ such that for every $D \ge 1$ and every polynomial
$U(z) = \sum_{n=1}^D a_n z^n$ with $U(0) = 0$, we have
\[
  \norm{U}_{\mathcal{F}} \le C\,\norm{\rho(U)}_\mathcal{F}.
\]
Here, $\rho(w)\coloneq ||1+w|-1|$ and one can take $C = 4600$.
\end{theorem}
\begin{remark}
  The condition $U(0) = 0$ means that $U$ is orthogonal to the constant
function $\mathbf{1}$ in $\mathcal{F}^2(\mathbb{C})$. Thus, the orthogonal reduction tells us that Theorem~\ref{thm:orthog} is strong enough to recover our main results.
\end{remark}
 %Since the passage from
%Theorem~\ref{thm:orthog} to Theorem~\ref{thm:local} will follow from a general
%Hilbert-space argument, an analogous reduction also holds in the general setting of Theorem \ref{thm:local_Gen} and its higher dimensional variants.

We now outline the steps needed to prove Theorem~\ref{thm:orthog}.
\medskip

\noindent
\textbf{Polar coordinates.}\;
A simple but important idea in our proof below is to represent the polynomial $U$ in Theorem \ref{thm:orthog} in polar coordinates. This will enable us to rephrase many of our estimates in terms of estimates for trigonometric polynomials (where we have more direct access to Fourier analytic tools). We further note that the Gaussian integral decomposes as
\begin{equation}\label{eq:polar}
\frac{1}{\pi}\int_{\mathbb{C}} F(z)\, \mathrm{e}^{-\lvert z \rvert^2}\,\mathrm{d} m(z)
\;=\;
2\int_0^\infty r\,\mathrm{e}^{-r^2}
\biggl(\int_{S^1} F(r\mathrm{e}^{\mathrm{i} t})\,\mathrm{d}\sigma(t)\biggr)\mathrm{d} r,
\end{equation}
where $\mathrm{d}\sigma = \mathrm{d} t/(2\pi)$ is the probability Haar measure on the
circle. Moreover, both norms in Theorem~\ref{thm:orthog} decompose in this way, so it
suffices to compare the inner circle integrals for each radius~$r$. This simple observation will be very useful in the forthcoming analysis.

\medskip
\noindent
\textbf{Circle estimates (Section~\ref{subsec:circle}).}\;
On a circle of radius~$r$, the function $t \mapsto U(r\mathrm{e}^{\mathrm{i} t})$ is a
trigonometric polynomial with strictly positive frequencies (since
$U(0) = 0$).  We prove two estimates relating the $L^2$ norm of such a
polynomial~$P$ to that of $\rho \circ P$:
\begin{enumerate}[(i)]
\item A \emph{local circle estimate} that works for any finite set of
  positive frequencies, at the cost of a constant depending on how many
  frequencies are present.

\item A \emph{high-frequency estimate} that applies when the lowest
  frequency is much larger than the bandwidth, giving a universal constant
  independent of the number of frequencies.
\end{enumerate}

\medskip
\noindent
\textbf{Frequency localization (Section~\ref{subsec:blocks}).}\; The polynomial $U$ may have up to $D$ frequencies, but the monomial~$z^n$
concentrates its Gaussian mass near the annulus $\lvert z \rvert \approx \sqrt{n}$.
We exploit this by partitioning the frequencies into blocks, each localized
to a specific annulus.  For each annulus, only a bounded window of nearby
blocks contributes significantly; the energy leaking from distant blocks
decays as a Gaussian in the block distance.

\medskip
\noindent
\textbf{Assembly (Section~\ref{subsec:assembly}).}\;
Here, we assemble the estimates of the previous two subsections to complete the proof of the orthogonal coercivity bound in Theorem \ref{thm:orthog}. The observation is that on each annulus, the nearby blocks form a trigonometric polynomial to which
one of the two circle estimates applies.  This gives a uniform bound on the
$\rho$-norm of the ``local'' polynomial.  We then use the $1$-Lipschitz
property of $\rho$ to pass from the local polynomial to the full $U$,
integrate over all annuli, and absorb the exponentially small leakage.

\subsection{The local reduction}\label{subsec:local-reduction}
In this subsection, we  reduce Theorem~\ref{main:local} to Theorem~\ref{thm:local}. The reduction is general. The first step is to show that local stable phase retrieval is implied by an
even more local variant. The second step is a simple density argument. The results in this section are the only results which have not been Lean verified, as they are not specific to STFT phase retrieval and are likely well-known.

\begin{proposition}[Local reduction]\label{prop:local-reduction}
Let \(H\) be a reproducing kernel Hilbert space of real or complex-valued
functions isometrically embedded in \(L^2(\Omega,\mu)\) where $(\Omega,\mu)$ is some measure space. Fix \(F\in H\) and
\(\epsilon>0\). Suppose that there exists a constant \(C>0\) such that for
all \(G\in H\) satisfying
\[
\inf_{\lvert \lambda \rvert=1}\lVert F-\lambda G \rVert_{L^2(\Omega)}<\epsilon
\]
we have
\[
\inf_{\lvert \lambda \rvert=1}\lVert F-\lambda G \rVert_{L^2(\Omega)}
\le
C\lVert \lvert F \rvert-\lvert G \rvert \rVert_{L^2(\Omega)}.
\]
Suppose also that phase retrieval holds at \(F\), i.e., if \(G\in H\) and
\(\lvert F \rvert=\lvert G \rvert\), then there exists \(\lambda\in\mathbb{C}\) with
\(\lvert \lambda \rvert=1\) and \(F=\lambda G\). Then there exists a constant
\(C'>0\) such that for all \(G\in H\),
\[
\inf_{\lvert \lambda \rvert=1}\lVert F-\lambda G \rVert_{L^2(\Omega)}
\le
C'\lVert \lvert F \rvert-\lvert G \rvert \rVert_{L^2(\Omega)}.
\]
\end{proposition}

\begin{proof}
We need to verify that there exists a constant \(C'\) such that for all
\(G\in H\),
\[
\inf_{\lvert \lambda \rvert=1}\lVert F-\lambda G \rVert_{L^2(\Omega)}
\le
C'\lVert \lvert F \rvert-\lvert G \rvert \rVert_{L^2(\Omega)}.
\]
First suppose that \(\lVert G \rVert_{L^2(\Omega)}\ge4\lVert F \rVert_{L^2(\Omega)}\).
Then
\[
\lVert \lvert F \rvert-\lvert G \rvert \rVert_{L^2(\Omega)}
\ge
\lVert G \rVert_{L^2(\Omega)}-\lVert F \rVert_{L^2(\Omega)}
\ge
\frac34\lVert G \rVert_{L^2(\Omega)},
\]
while
\[
\inf_{\lvert \lambda \rvert=1}\lVert F-\lambda G \rVert_{L^2(\Omega)}
\le
\lVert F \rVert_{L^2(\Omega)}+\lVert G \rVert_{L^2(\Omega)}
\le
\frac54\lVert G \rVert_{L^2(\Omega)}.
\]
Thus the desired inequality holds in this case with any \(C'\ge5/3\).
Therefore, we may restrict attention to \(G\) with
\(\lVert G \rVert_{L^2(\Omega)}<4\lVert F \rVert_{L^2(\Omega)}\).
\medskip


Now suppose that the desired inequality is true for all \(G\) satisfying
\[
\inf_{\lvert \lambda \rvert=1}\lVert F-\lambda G \rVert_{L^2(\Omega)}<\epsilon,
\]
but is not true in general. Then for every \(n\) we can find \(G_n\in H\)
with
\[
\inf_{\lvert \lambda \rvert=1}\lVert F-\lambda G_n \rVert_{L^2(\Omega)}\ge\epsilon,
\]
\[
\inf_{\lvert \lambda \rvert=1}\lVert F-\lambda G_n \rVert_{L^2(\Omega)}
>
n\lVert \lvert F \rvert-\lvert G_n \rvert \rVert_{L^2(\Omega)},
\]
and also \(\lVert G_n \rVert_{L^2(\Omega)}<4\lVert F \rVert_{L^2(\Omega)}\). This latter
inequality guarantees that
\(\inf_{\lvert \lambda \rvert=1}\lVert F-\lambda G_n \rVert_{L^2(\Omega)}\) is uniformly
bounded, and hence
\[
\lVert \lvert F \rvert-\lvert G_n \rvert \rVert_{L^2(\Omega)}\to0.
\]
In particular, after passing to a subsequence, we may assume that
\[
\lvert G_{n_k}(x) \rvert\to\lvert F(x) \rvert
\qquad\text{for almost every }x\in\Omega.
\]
Since \(\lVert G_{n_k} \rVert_{L^2(\Omega)}\) is bounded independently of \(k\) and
the embedding of \(H\) into \(L^2(\Omega,\mu)\) is isometric, the sequence
\((G_{n_k})\) is bounded in \(H\). Passing to a further subsequence if
necessary, we may assume that \(G_{n_{k_\ell}}\rightharpoonup G\) weakly in
\(H\). Since point evaluations are bounded on \(H\), we have
\[
G_{n_{k_\ell}}(x)\to G(x)
\]
for every \(x\in\Omega\). Hence \(\lvert G_{n_{k_\ell}}(x) \rvert\to\lvert G(x) \rvert\) for
every \(x\in\Omega\). This forces \(\lvert G \rvert=\lvert F \rvert\) almost everywhere. By the
phase-retrieval hypothesis, \(G=\lambda F\) for some unimodular \(\lambda\).
\medskip

Collecting what we know, \(G_{n_{k_\ell}}\) converges weakly in \(H\) to
\(\lambda F\). Moreover, from
\(\lVert \lvert F \rvert-\lvert G_n \rvert \rVert_{L^2(\Omega)}\to0\), we see that
\[
\lVert G_n \rVert_{L^2(\Omega)}=\lVert \lvert G_n \rvert \rVert_{L^2(\Omega)}
\to
\lVert \lvert F \rvert \rVert_{L^2(\Omega)}
=\lVert \lambda F \rVert_{L^2(\Omega)}.
\]
Because the embedding is isometric, weak convergence in \(H\) together with
convergence of the \(L^2\)-norms to the norm of the weak limit implies strong
convergence in \(L^2\). Thus \(G_{n_{k_\ell}}\to\lambda F\) in
\(L^2(\Omega)\). For all sufficiently large \(\ell\), this contradicts
\[
\inf_{\lvert \lambda \rvert=1}
\lVert F-\lambda G_{n_{k_\ell}} \rVert_{L^2(\Omega)}
\ge \epsilon.
\]
\end{proof}
\subsubsection{Proof of Theorem \ref{main:local}} Using this reduction, we immediately deduce Theorem \ref{main:local} from Theorem \ref{thm:local} from density and by applying Proposition \ref{prop:local-reduction} with $F=1$ and $H=\mathcal{F}^2(\mathbb{C})$.
\subsection{The orthogonal reduction}\label{subsec:reduction}
In this subsection, we work in the generality of a complex Hilbert space $H$
isometrically embedded in $L^2(\mu)$ for some measure space $(\Omega, \mu)$.
We show that \emph{orthogonal coercivity} -- a stability estimate for
perturbations perpendicular to a fixed unit vector -- implies \emph{local
stability} for all sufficiently small perturbations, after optimising over a
global phase.  The argument is purely Hilbert-space-theoretic and uses no
structure specific to the Fock space. 

\begin{theorem}[Orthogonal reduction]\label{thm:reduction}
\par
Let $H$ be a complex Hilbert space with norm $\lVert \cdot \rVert$, isometrically
embedded in $L^2(\mu)$.  Let $f_0 \in H$ with $\lVert f_0 \rVert = 1$ and
suppose that there exists $M > 0$ such that
\begin{equation}\label{eq:orthog-hyp}
\lVert g \rVert
\;\le\;
M\, \bigl\|\, \lvert f_0 + g \rvert - \lvert f_0 \rvert \,\bigr\|_{L^2}
\qquad
\text{for all } g \perp f_0.
\end{equation}
Then there exist $\delta > 0$ and $\widetilde{M} > 0$, depending only
on~$M$, such that
for every $h \in H$ with $\lVert h \rVert \le \delta$ and
$\langle h,f_0\rangle \in \mathbb{R}$, we have
\[
\lVert h \rVert
\;\le\;
\widetilde{M}\, \bigl\|\, \lvert f_0 + h \rvert - \lvert f_0 \rvert \,\bigr\|_{L^2}.
\]
One may take $\delta = 1/(M+1)$ and $\widetilde{M} = 5M + 3$.
\end{theorem}

\begin{remark}
The condition $\langle h,f_0\rangle \in \mathbb{R}$ is a phase normalisation: it says that
$h$ has been rotated so that its component along $f_0$ is real.  In the
application to Theorem~\ref{thm:local}, this is achieved by choosing the
optimal phase $w$ so that $\langle w(f_0 + p) - f_0,f_0\rangle$ is real and
non-negative, which amounts to aligning the phase of the inner product
$\langle p,f_0\rangle$.
\end{remark}

\begin{proof}
Write $m = \bigl\|\, \lvert f_0 + h \rvert - \lvert f_0 \rvert \,\bigr\|_{L^2}$ for the right-hand side of our inequality.

\medskip
\noindent
\emph{Step 1: Orthogonal decomposition.}\;
Set $a = \operatorname{Re}\langle h,f_0\rangle = \langle h,f_0\rangle$ (a real number by hypothesis) and
$g = h - a f_0$.  Then $g \perp f_0$, $h = af_0 + g$, and
$\lVert h \rVert^2 = a^2 + \lVert g \rVert^2$, so
\begin{equation}\label{eq:triangle}
\lVert h \rVert \;\le\; \lvert a \rvert + \lVert g \rVert.
\end{equation}

\medskip
\noindent
\emph{Step 2: Control of the scalar coefficient.}\;
The pointwise identity
$\lvert f_0 + h \rvert^2 - \lvert f_0 \rvert^2
= 2\,\operatorname{Re}(h\,\overline{f_0}) + \lvert h \rvert^2$
integrates to
\[
\int_\Omega \bigl(\lvert f_0 + h \rvert^2 - \lvert f_0 \rvert^2\bigr)\,\mathrm{d}\mu
= 2a + \lVert h \rVert^2.
\]
On the other hand, the factorization
$\lvert f_0 + h \rvert^2 - \lvert f_0 \rvert^2
= \bigl(\lvert f_0 + h \rvert - \lvert f_0 \rvert\bigr)
  \bigl(\lvert f_0 + h \rvert + \lvert f_0 \rvert\bigr)$
and Cauchy--Schwarz gives
\[
\lvert 2a + \lVert h \rVert^2 \rvert
\;\le\;
m \cdot \bigl\|\, \lvert f_0 + h \rvert + \lvert f_0 \rvert \,\bigr\|_{L^2}
\;\le\;
m\,(2 + \lVert h \rVert).
\]
Hence,
\begin{equation}\label{eq:a-bound}
\lvert a \rvert
\;\le\;
\frac{2 + \lVert h \rVert}{2}\, m
+ \frac{\lVert h \rVert^2}{2}.
\end{equation}

\medskip
\noindent
\emph{Step 3: Control of the orthogonal component.}\;
The reverse triangle inequality
$\bigl\lvert \lvert f_0 + g \rvert - \lvert f_0 + h \rvert \bigr\rvert
\le \lvert g - h \rvert = \lvert a \rvert\,\lvert f_0 \rvert$
gives
$\bigl\|\, \lvert f_0 + g \rvert - \lvert f_0 \rvert \,\bigr\|_{L^2}
\le \lvert a \rvert + m$.
The orthogonal coercivity hypothesis~\eqref{eq:orthog-hyp} then yields
\begin{equation}\label{eq:g-bound}
\lVert g \rVert \;\le\; M\,(\lvert a \rvert + m).
\end{equation}

\medskip
\noindent
\emph{Step 4: Completing the proof.}\;
Substituting~\eqref{eq:a-bound} and~\eqref{eq:g-bound}
into~\eqref{eq:triangle}:
\[
\lVert h \rVert
\;\le\;
\lvert a \rvert + M(\lvert a \rvert + m)
= (M+1)\lvert a \rvert + Mm
\;\le\;
\underbrace{\Bigl(M + \frac{(M+1)(2+\lVert h \rVert)}{2}\Bigr)}_{=:\, C(\lVert h \rVert)}
m
\;+\; \frac{M+1}{2}\,\lVert h \rVert^2.
\]
Since $\lVert h \rVert \le \delta$, we have $\lVert h \rVert^2 \le \delta\,\lVert h \rVert$, so
\[
\lVert h \rVert
\;\le\; C(\delta)\, m + \frac{(M+1)\delta}{2}\,\lVert h \rVert.
\]
Choosing $\delta = 1/(M+1)$ makes $(M+1)\delta/2 = 1/2$, giving
$\lVert h \rVert/2 \le C(\delta)\, m$ and therefore
\[
\lVert h \rVert \;\le\; 2\, C(\delta)\, m.
\]
A short computation shows that
$C(\delta) = M + (M+1)(2 + \delta)/2 < (5M+3)/2$, so one may take
$\widetilde{M} = 5M + 3$.
\end{proof}

\begin{corollary}[Fock-space application]\label{cor:fock-local}
Taking $f_0 = \mathbf{1}$ (the constant function), $H = \mathcal{F}^2(\mathbb{C})$, and
$M = 4600$ (from Theorem~\ref{thm:orthog}), the orthogonal reduction gives
$\delta = 1/4601$ and $\widetilde{M} = 23003$.
This is Theorem~\ref{thm:local}.
\end{corollary}

\begin{proof}
The condition $g \perp \mathbf{1}$ in $\mathcal{F}^2(\mathbb{C})$ means precisely that
$g(0) = 0$, i.e., $g$ has no constant term.
Theorem~\ref{thm:orthog} provides the orthogonal coercivity
hypothesis~\eqref{eq:orthog-hyp} with $M = 4600$.
The optimal phase $w$ is chosen so that
$h = w(1 + p) - 1$ satisfies $\langle h,\mathbf{1}\rangle \in \mathbb{R}$ and
$\lVert h \rVert_\mathcal{F} \le \delta$; the existence of such a $w$ is straightforward.
Theorem~\ref{thm:reduction} then yields the conclusion of
Theorem~\ref{thm:local}. 
\end{proof}
Our goal for the remaining subsections will  be to establish the orthogonal coercivity bound in Theorem  \ref{thm:orthog}.
 \subsection{Circle estimates}\label{subsec:circle}

In this section we prove two estimates controlling the $L^2$ norm of a
trigonometric polynomial $P$ on the circle in terms of the $L^2$ norm of
$\rho \circ P$.  Both exploit the fact that all frequencies of~$P$ are
strictly positive (which is a consequence of our reduction to the case $U(0)=0$) but they apply in
complementary regimes.  The \emph{local circle estimate}
(Theorem~\ref{thm:local-circle}) works for any frequency set and gives a
constant proportional to $\sqrt{L}$, where $L$ is the number of frequencies.
The \emph{high-frequency estimate} (Theorem~\ref{thm:high-freq}) applies
when the lowest frequency is much larger than the bandwidth, and gives an
absolute constant. Together they will provide a uniform bound on every
annulus.
\medskip

Throughout this section, $P$ denotes a trigonometric polynomial
$P(t) = \sum_{n \in E} b_n \mathrm{e}^{\mathrm{i} nt}$ for a finite set $E \subset \mathbb{Z}_{>0}$ 
and all norms are taken in $L^2(S^1, \mathrm{d}\sigma)$. Recall as well that the function $\rho(w) = \bigl\lvert \lvert 1+w \rvert - 1 \bigr\rvert$
 is $1$-Lipschitz, and in particular we have $\rho(w) \le \lvert w \rvert$ and
 %by two applications of the reverse triangle inequality:
%$\lvert \rho(w) - \rho(z) \rvert
%\le \bigl\lvert \lvert 1+w \rvert - \lvert 1+z \rvert \bigr\rvert
%\le \lvert w - z \rvert$.
%In particular, $\rho(w) \le \lvert w \rvert$ (take $z = 0$) and
\begin{equation}\label{eq:rho-lip}
\rho(w) \;\ge\; \rho(z) - \lvert w - z \rvert.
\end{equation}

\subsubsection{The local circle estimate}\label{sec:local}

The following elementary inequality will be  used several times in the forthcoming calculations, so we record it here for convenience.

\begin{lemma}\label{lem:safe-square}
Let $a, b, c \ge 0$ with $a \ge b - c$.  Then
$a^2 \ge \tfrac{1}{2}\, b^2 - c^2$.
\end{lemma}

\begin{proof}
If $b \le c$ the right hand-side is nonpositive.  If $b > c$ then
$a \ge b-c > 0$ gives $a^2 \ge (b-c)^2 \ge \tfrac{1}{2}b^2 - c^2$, where
the last step is the identity
$(b-c)^2 - \tfrac{1}{2}b^2 + c^2 = \tfrac{1}{2}(b-2c)^2 \ge 0$.
\end{proof}

\begin{theorem}[Local circle estimate]\label{thm:local-circle} Let $E \subset \mathbb{Z}_{>0}$ with $\lvert E \rvert = L \ge 1$, let $b_n \in \mathbb{C}$ for
$n \in E$, and set $P(t) = \sum_{n \in E} b_n\,\mathrm{e}^{\mathrm{i} nt}$.  Then
\[
\lVert P \rVert_{L^2}^2
\;\le\;
144\, L \cdot \lVert \rho \circ P \rVert_{L^2}^2.
\]
\end{theorem}

\begin{proof}
Write $\alpha = \lVert P \rVert_{L^2}$ and $\beta = \lVert \rho \circ P \rVert_{L^2}$.
We show $\beta \ge \alpha/(12\sqrt{L})$, which gives the result after
squaring. To prove this estimate, we will consider two different regimes; namely, the one where $\alpha$ is sufficiently small and its complement. The threshold between the two regimes will be taken to be
$\alpha=\alpha_0 = 1/(4\sqrt{L})$.
\medskip

Two structural properties of $P$ underpin both cases.  First,
by Cauchy--Schwarz and Parseval we have the pointwise bound
$\lvert P(t) \rvert \le \sqrt{L}\, \alpha$ for all~$t$.
Second, since every $n \in E$ satisfies $n \ge 1$, the zeroth Fourier
coefficient of $P$ vanishes, i.e.,
$\int_{S^1} P\,\mathrm{d}\sigma = 0$.
More importantly, we have $n + m \ge 2$ for all $n, m \in E$, so
$\int_{S^1} P^2\,\mathrm{d}\sigma = 0$ as well.  Writing $P = u + \mathrm{i} v$ with
$u, v$ real-valued, this identity splits as
\begin{equation}\label{eq:equal-mass}
\int_{S^1} u^2\,\mathrm{d}\sigma
= \int_{S^1} v^2\,\mathrm{d}\sigma
= \frac{\alpha^2}{2}
\qquad\text{and}\qquad
\int_{S^1} uv\,\mathrm{d}\sigma = 0.
\end{equation}
This equipartition of mass between the real and imaginary parts of $P$ will be crucial in the small amplitude regime below.
\medskip

\noindent
\textbf{Small amplitudes} ($\alpha \le \alpha_0$).
Here, we have $\lVert P \rVert_\infty \le \sqrt{L}\,\alpha \le 1/4$.  For $\lvert w \rvert \le 1/4$
we also have $\lvert 1+w \rvert \ge 3/4$ and $\operatorname{Re}(w) \ge -1/4$, so the denominator
$\lvert 1+w \rvert + 1 + \operatorname{Re}(w) \ge 3/2 > 0$.  The algebraic identity
\[
\lvert 1+w \rvert - 1 - \operatorname{Re}(w)
= \frac{\operatorname{Im}(w)^2}{\lvert 1+w \rvert + 1 + \operatorname{Re}(w)}
\]
shows that $\lvert 1+w \rvert - 1 = \operatorname{Re}(w) + \varepsilon$ with
$0 \le \varepsilon \le \operatorname{Im}(w)^2/(3/2) \le \lvert w \rvert^2$.  In particular,
$\rho(w) = \bigl\lvert \operatorname{Re}(w) + \varepsilon\bigr\rvert \ge
\lvert \operatorname{Re}(w) \rvert - \varepsilon$, so by 
Lemma~\ref{lem:safe-square} we have
\[
\rho(P)^2 \;\ge\; \tfrac{1}{2}\, u^2 - \lvert P \rvert^4.
\]
In other words, when $P$ is small in amplitude, most of the mass of $\rho(P)$ comes from the real part of $P$. One might a priori worry about the hypothetical possibility that if $P(t)$ were allowed to be purely imaginary for most $t$, the nonlinear function $\rho(P(t))$ might be small relative to $P(t)$. However, integrating and using the equipartition property ~\eqref{eq:equal-mass} ensures that we do indeed have the coercive bound:
$\beta^2 \ge \alpha^2/4 - \lVert P \rVert_\infty^2\, \alpha^2 \ge \alpha^2/4 - L\,\alpha^4$.
Since $L\alpha^2 \le 1/16$, we get $\beta^2 \ge 3\alpha^2/16$, so
$\beta \ge \sqrt{3}\,\alpha/4 \ge \alpha/(12\sqrt{L})$.

\medskip
\noindent
\textbf{Large amplitudes} ($\alpha > \alpha_0$).
The fact that $P$ has mean zero (since it is supported on positive frequencies) also gives
$\int_{S^1}\lvert 1+P \rvert^2\,\mathrm{d}\sigma = 1 + \alpha^2$, so
\[
\alpha^2
= \int_{S^1}\bigl(\lvert 1+P \rvert^2 - 1\bigr)\,\mathrm{d}\sigma
\le \int_{S^1} \rho(P)\,\bigl(\lvert 1+P \rvert + 1\bigr)\,\mathrm{d}\sigma,
\]
where we used $\bigl\lvert \lvert 1+P \rvert^2 - 1\bigr\rvert
= \bigl(\lvert 1+P \rvert + 1\bigr)\rho(P)$.  By Cauchy--Schwarz,
\[
\alpha^2 \;\le\; \beta \cdot \bigl(\sqrt{1+\alpha^2} + 1\bigr).
\]
If $\alpha \le 1$, then $\beta \ge \alpha^2/3$ and $\alpha > \alpha_0 = 1/(4\sqrt{L})$ give
$\beta \ge \alpha/(12\sqrt{L})$.
If $\alpha > 1$, then $\beta \ge \alpha/3 \ge \alpha/(12\sqrt{L})$.
In both regimes, we have $\beta \ge \alpha/(12\sqrt{L})$.
\end{proof}


\subsubsection{Rotational averaging}\label{sec:rot-avg}

The local circle estimate grows with the number of frequencies~$L$.  To
eliminate this dependency for high-frequency polynomials, we will need the
following fact: the $\rho$-image of a full circle already captures a
definite fraction of its radius.

\begin{proposition}[Rotational averaging]\label{prop:rot-avg}
For all $r \ge 0$, we have
\[
\int_{S^1} \rho(r\mathrm{e}^{\mathrm{i}\theta})^2\,\mathrm{d}\sigma(\theta)
\;\ge\; \frac{r^2}{8}.
\]
\end{proposition}

\begin{proof}
On the arc $\lvert \theta \rvert \le \pi/4$ (which has $\mathrm{d}\sigma$-measure $1/4$),
we claim that $\rho(r\mathrm{e}^{\mathrm{i}\theta}) \ge r/\sqrt{2}$.  Indeed,
$\cos\theta \ge 1/\sqrt{2}$ gives
\[
\lvert 1 + r\mathrm{e}^{\mathrm{i}\theta} \rvert - 1
= \frac{r^2 + 2r\cos\theta}{\lvert 1+r\mathrm{e}^{\mathrm{i}\theta} \rvert + 1}
\ge \frac{r(r + \sqrt{2})}{r + 2}
\ge \frac{r}{\sqrt{2}}.
\]
  Moreover, the numerator
is positive, so $\lvert 1 + r\mathrm{e}^{\mathrm{i}\theta} \rvert > 1$ and $\rho$ equals
$\lvert 1+r\mathrm{e}^{\mathrm{i}\theta} \rvert - 1$ without the absolute value.  Restricting the
integral to this arc, we have
$\int \rho^2\,\mathrm{d}\sigma \ge \frac{1}{4}(r/\sqrt{2})^2 = r^2/8$.
\end{proof}


\subsubsection{The high-frequency estimate}\label{sec:high-freq}

When the polynomial oscillates rapidly, the Lipschitz property
of~$\rho$ and the rotational averaging of
Proposition~\ref{prop:rot-avg} combine to give an absolute bound.  The idea
is to factor $P(t) = \mathrm{e}^{\mathrm{i} Nt}\, Q(t)$ where $Q$ is a slowly varying
envelope. On short intervals, the leading phase $\mathrm{e}^{\mathrm{i} Nt}$ sweeps through
a full rotation while $Q$ barely changes, and the rotational averaging
ensures that $\rho$ captures a positive fraction of $\lvert Q \rvert^2$.  The only
analytic input is a Poincar\'e inequality quantifying how well $Q$ is
approximated by its local averages.

\begin{theorem}[High-frequency estimate]\label{thm:high-freq} Let $N \ge 1$ and $L \ge 1$ satisfy $1343\, L^2 \le N^2$.
For any $b_0, \dots, b_{L-1} \in \mathbb{C}$, the polynomial
$P(t) = \sum_{m=0}^{L-1} b_m\, \mathrm{e}^{\mathrm{i}(N+m)t}$ satisfies
\[
\lVert P \rVert_{L^2}^2
\;\le\;
32\, \lVert \rho \circ P \rVert_{L^2}^2.
\]
\end{theorem}

\begin{proof}
Write $P(t) = \mathrm{e}^{\mathrm{i} Nt}\,Q(t)$ with
$Q(t) = \sum_{m=0}^{L-1} b_m\,\mathrm{e}^{\mathrm{i} mt}$, so $\lVert P \rVert = \lVert Q \rVert$.
Partition the circle into $N$ equal arcs
$J_k = [2\pi k/N,\; 2\pi(k+1)/N)$, and let
$q_k = N\!\int_{J_k}\! Q\,\mathrm{d}\sigma$ be the average of $Q$ on~$J_k$.
\medskip

On each interval $J_k$ of length $h = 2\pi/N$, the Poincar\'e
inequality gives
$\int_{J_k}\lvert Q - q_k \rvert^2\,\mathrm{d}\sigma \le h^2\int_{J_k}\lvert Q' \rvert^2\,
\mathrm{d}\sigma$.\footnote{We use the inequality
$\int_0^h \lvert f - \bar f \rvert^2\,\mathrm{d} x
\le h^2\int_0^h\lvert f' \rvert^2\,\mathrm{d} x$
with the non-sharp constant~$1$; here $\bar f$ denotes the mean of~$f$
on~$[0,h]$.  The sharp 
constant is $1/\pi^2$, which, amusingly, we saw after the conversion of the lean code to natural language that the LLM emphasized would improve the threshold from $1343$ to~$136$ but did not use it as it is not yet in MathLib.}
Summing over $k$ and using Bernstein's inequality $\lVert Q' \rVert^2 \le L^2\lVert Q \rVert^2$
(since every frequency of~$Q$ has modulus at most~$L{-}1 < L$) we see that
\begin{equation}\label{eq:delta}
\delta \;:=\; \sum_{k} \int_{J_k}\lvert Q - q_k \rvert^2\,\mathrm{d}\sigma
\;\le\; \frac{4\pi^2 L^2}{N^2}\,\lVert Q \rVert^2.
\end{equation}
Note that
\begin{equation}\label{eq:parseval-partition}
\lVert Q \rVert^2
= \frac{1}{N}\sum_{k}\lvert q_k \rvert^2 + \delta.
\end{equation}
As $t$ traverses $J_k$, the argument $Nt$ traverses one full period.
Substituting $s = Nt$ and applying Proposition~\ref{prop:rot-avg} with
$r = \lvert q_k \rvert$ we obtain
\[
\int_{J_k}\rho\bigl(\mathrm{e}^{\mathrm{i} Nt} q_k\bigr)^2\,\mathrm{d}\sigma
= \frac{1}{N}\int_{S^1}\rho\bigl(\lvert q_k \rvert\mathrm{e}^{\mathrm{i} s}\bigr)^2\,\mathrm{d}\sigma(s)
\;\ge\; \frac{\lvert q_k \rvert^2}{8N}.
\]
On $J_k$, the Lipschitz bound~\eqref{eq:rho-lip} gives
$\rho(P) = \rho(\mathrm{e}^{\mathrm{i} Nt}Q)
\ge \rho(\mathrm{e}^{\mathrm{i} Nt}q_k) - \lvert Q - q_k \rvert$.
By Lemma~\ref{lem:safe-square},
$\rho(P)^2 \ge \tfrac{1}{2}\rho(\mathrm{e}^{\mathrm{i} Nt}q_k)^2 - \lvert Q - q_k \rvert^2$.
Integrating over $J_k$ and using the above we obtain
\[
\int_{J_k}\rho(P)^2\,\mathrm{d}\sigma
\;\ge\; \frac{\lvert q_k \rvert^2}{16N} - \int_{J_k}\lvert Q-q_k \rvert^2\,\mathrm{d}\sigma.
\]
Summing over $k$ and applying~\eqref{eq:parseval-partition}, it follows that
\[
\lVert \rho \circ P \rVert^2
\;\ge\; \frac{1}{16}\Bigl(\lVert Q \rVert^2 - \delta\Bigr) - \delta
= \frac{\lVert Q \rVert^2}{16} - \frac{17\delta}{16}.
\]
Substituting the bound~\eqref{eq:delta} we arrive at
\[
\lVert \rho \circ P \rVert^2
\;\ge\;
\frac{1}{16}\Bigl(1 - \frac{68\pi^2 L^2}{N^2}\Bigr)\lVert Q \rVert^2.
\]
The condition $1343\,L^2 \le N^2$ gives $68\pi^2 L^2/N^2 \le 1/2$
(since $136\pi^2 < 1343$), and therefore
$\lVert \rho \circ P \rVert^2 \ge \lVert P \rVert^2/32$.
\end{proof}

\subsection{Frequency localization}\label{subsec:blocks}
The circle estimates from Section~\ref{subsec:circle} require the trigonometric
polynomial to have a controlled number of frequencies, yet $U$ may have
arbitrarily many.  The Gaussian weight is what  resolves this issue: the radial factor
$r \mapsto r^{2n+1}\mathrm{e}^{-r^2}$ peaks sharply near $r = \sqrt{n + 1/2}$,
so grouping frequencies by peak location decomposes $U$ into ``blocks''
approximately supported on disjoint annuli.
\medskip

We begin with some general notation. For $\ell \ge 1$, we define the \emph{frequency block}
\[
I_\ell = \bigl\{n \in \mathbb{N} : \ell^2 \le n < (\ell+1)^2\bigr\},
\]
which has $\lvert I_\ell \rvert = 2\ell+1$ elements, and the corresponding
\emph{block polynomial}
$U_\ell(z) = \sum_{n \in I_\ell,\, n \le D} a_n z^n$. The blocks $I_1, I_2, \dots$ are pairwise disjoint and cover
$\{1, \dots, D\}$ up to index $\Lambda = \lfloor\sqrt{D}\rfloor$.  Since the
Fock norm is a sum over individual frequencies, it decomposes as
\begin{equation}\label{eq:fock-decomp}
\lVert U \rVert_\mathcal{F}^2 = \sum_{\ell=1}^{\Lambda} \lVert U_\ell \rVert_\mathcal{F}^2.
\end{equation}
Similarly, on each circle of radius~$r$, blocks with distinct frequency
supports are orthogonal, so we have
$\int_{S^1} U_{\ell_1}(r\mathrm{e}^{\mathrm{i} t})\,
\overline{U_{\ell_2}(r\mathrm{e}^{\mathrm{i} t})}\,\mathrm{d}\sigma(t) = 0$
for $\ell_1 \ne \ell_2$.  We record the following obvious lemma to set notation.

\begin{lemma}\label{lem:peak}
For $n \in I_\ell$ with $\ell \ge 1$, 
$r_n = \sqrt{n + 1/2}$ satisfies $\ell < r_n < \ell + 1$.
\end{lemma}

Now fix a window width $M \ge 1$.  For each annulus index $j \ge 0$, define the
\emph{local polynomial}
\[
V_j = \sum_{\substack{1 \le \ell \le \Lambda \\ \lvert \ell - j \rvert \le M}} U_\ell
\qquad\text{and the \emph{remainder}}\qquad
R_j = U - V_j.
\]
The local polynomial gathers the blocks whose peaks fall within $M$ steps of
annulus~$j$. For $j \ge M + 1$, the frequency support of $V_j$ lies in
$[(j{-}M)^2,\; (j{+}M{+}1)^2)$, and the number of frequencies is
$L_j = (2M{+}1)(2j{+}1)$.
\medskip


The energy that a block $U_\ell$ contributes to a distant annulus decays
like $\mathrm{e}^{-(\lvert j-\ell \rvert-1)^2}$.  The mechanism is simple: the function
$\varphi_n(r) = (2n{+}1)\log r - r^2$ (so that
$r^{2n+1}\mathrm{e}^{-r^2} = \mathrm{e}^{\varphi_n(r)}$) has a unique maximum at
$r_n = \sqrt{n+1/2}$ and concavity $\varphi_n''(r) \le -2$, giving the
Gaussian upper bound
\begin{equation}\label{eq:phi-quad}
\mathrm{e}^{\varphi_n(r)}
\;\le\;
\mathrm{e}^{\varphi_n(r_n)}\, \mathrm{e}^{-(r - r_n)^2}.
\end{equation}
The value at the saddle point is controlled as follows.

\begin{lemma}\label{lem:exp-phi}
\par\noindent
For $n \ge 1$,\;
$\mathrm{e}^{\varphi_n(r_n)} \le \mathrm{e}^{1/4}\, n!/2$.
\end{lemma}

\begin{proof}
The Gamma integral gives
$n!/2 = \int_0^\infty \mathrm{e}^{\varphi_n(r)}\,\mathrm{d} r$.
Since $r_n \ge \sqrt{3/2} > 1$, the interval $[r_n - 1/2,\, r_n + 1/2]$
lies in $(0,\infty)$, and~\eqref{eq:phi-quad} shows that
$\varphi_n \ge \varphi_n(r_n) - 1/4$ on this interval.  Hence
$n!/2 \ge \mathrm{e}^{\varphi_n(r_n) - 1/4}$.
\end{proof}

Combining~\eqref{eq:phi-quad} and Lemma~\ref{lem:exp-phi}, we see that for $r \in
[j, j{+}1]$ and $n \in I_\ell$, the monomial weight satisfies
\begin{equation}\label{eq:monomial-bound}
r^{2n+1}\mathrm{e}^{-r^2}
\;\le\;
\frac{\mathrm{e}^{1/4}\, n!}{2}\,
\mathrm{e}^{-\operatorname{dist}(r_n,\,[j,j{+}1])^2},
\end{equation}
where $\operatorname{dist}(r_n, [j,j{+}1]) = \max(j - r_n,\, r_n - (j{+}1),\, 0)$.
By Lemma~\ref{lem:peak}, $r_n \in (\ell, \ell{+}1)$, so this distance is at
least $(\lvert j - \ell \rvert - 1)_+$. We now quantify the ``leakage" between the blocks.

\begin{proposition}\label{prop:single-leak}

For $\ell \ge 1$ and $j \ge 0$, we have
\[
\frac{1}{\pi}\int_{\{j \le \lvert z \rvert < j+1\}}
\lvert U_\ell(z) \rvert^2\, \mathrm{e}^{-\lvert z \rvert^2}\,\mathrm{d} m(z)
\;\le\;
\mathrm{e}^{1/4}\,
\mathrm{e}^{-(\lvert j-\ell \rvert-1)_+^2}\,
\lVert U_\ell \rVert_\mathcal{F}^2.
\]
\end{proposition}

\begin{proof}
Passing to polar coordinates and using the orthogonality of distinct
monomials on each circle, the left hand-side equals
$2\sum_{n \in I_\ell}\lvert a_n \rvert^2 \int_j^{j+1} r^{2n+1}\mathrm{e}^{-r^2}\mathrm{d} r$.
Bounding each radial integral by~\eqref{eq:monomial-bound} and summing gives the result.
\end{proof}

Summing over all distant blocks and swapping the order of summation
yields the following result.

\begin{proposition}\label{prop:total-leak} For $M \ge 2$, define $\eta_M = 2\mathrm{e}^{1/4}\sum_{m \ge M}\mathrm{e}^{-m^2}$.  Then
\[
\sum_{j \ge 0}\;
\frac{1}{\pi}\int_{\{j \le \lvert z \rvert < j+1\}}
\lvert R_j(z) \rvert^2\, \mathrm{e}^{-\lvert z \rvert^2}\,\mathrm{d} m(z)
\;\le\;
\eta_M \cdot \lVert U \rVert_\mathcal{F}^2.
\]
\end{proposition}

\begin{proof}
By the orthogonality on each annulus,
$\int \lvert R_j \rvert^2 = \sum_{\lvert \ell-j \rvert>M}\int\lvert U_\ell \rvert^2$.
Now apply Proposition~\ref{prop:single-leak}, swap $\sum_j$ and $\sum_\ell$
(all terms are nonneg\-ative), and note that for $\lvert j-\ell \rvert > M \ge 2$
the decay factor is $\mathrm{e}^{-(\lvert j-\ell \rvert-1)^2}$.  The inner sum over
$j$ contributes at most $2\sum_{m \ge M}\mathrm{e}^{-m^2}$ and multiplying by
$\mathrm{e}^{1/4}$ and using~\eqref{eq:fock-decomp} gives the result.
\end{proof}

%For $M = 5$, the geometric series bound
%$\sum_{m \ge 5}\mathrm{e}^{-m^2} \le \mathrm{e}^{-25}/(1 - \mathrm{e}^{-11}) < 1.39 \times
%10^{-11}$ gives the following result.

%\begin{corollary}\label{cor:eta5} $\eta_5 < 4 \times 10^{-11}$.
%\end{corollary}
We can now state the uniform estimate that holds on every annulus.  The
function $\zeta \mapsto V_j(r\zeta)$ is a trigonometric polynomial on $S^1$
with all frequencies in $\mathbb{Z}_{>0}$, and we can apply either
Theorem~\ref{thm:local-circle} or Theorem~\ref{thm:high-freq} depending
on~$j$.

\begin{proposition}\label{prop:annulus} With $M = 5$, we have for every $j \ge 0$ and every $r \in [j, j{+}1]$,
\[
\int_{S^1}\lvert V_j(r\zeta) \rvert^2\,\mathrm{d}\sigma(\zeta)
\;\le\;
1620^2 \cdot
\int_{S^1}\rho\bigl(V_j(r\zeta)\bigr)^2\,\mathrm{d}\sigma(\zeta).
\]
\end{proposition}

\begin{proof}
When $r = 0$ both sides vanish (all monomials have positive degree).
For $r > 0$, the function $\zeta \mapsto V_j(r\zeta)$ is a trigonometric
polynomial with positive frequencies and coefficients $a_n r^n$. We keep the amusing case distinction made during the autoformalization.
\medskip

\smallskip\noindent
\emph{Case $j \le 817$.}\;
With $M = 5$, the frequency count satisfies $L_j \le 11(2 \cdot 817 + 1)
< 18{,}000$.  By Theorem~\ref{thm:local-circle},
$\lVert V_j(r\,\cdot\,) \rVert^2 \le 144\, L_j \cdot
\lVert \rho \circ V_j(r\,\cdot\,) \rVert^2$, and
$144 \cdot 18{,}000 < 1620^2$.
\medskip

\smallskip\noindent
\emph{Case $j \ge 818$.}\;
The lowest frequency is $N_j = (j{-}5)^2$ and the bandwidth is
$L_j = 11(2j{+}1)$.  The ratio $L_j^2/N_j^2$ is decreasing in~$j$
(for $j \ge 6$), so it suffices to verify the condition
$1343\,L_j^2 \le N_j^2$ of Theorem~\ref{thm:high-freq} at $j = 818$.
A direct computation confirms $1343 \cdot (11 \cdot 1637)^2 < 813^4$.
The high-frequency estimate then gives
$\lVert V_j(r\,\cdot\,) \rVert^2 \le 32\,\lVert \rho \circ V_j(r\,\cdot\,) \rVert^2
\le 1620^2\,\lVert \rho \circ V_j(r\,\cdot\,) \rVert^2$.
\end{proof} 
\subsection{Proof of the orthogonal coercivity estimate}\label{subsec:assembly} With the estimates from subsections~\ref{subsec:circle} and \ref{subsec:blocks} and 
(Proposition~\ref{prop:annulus}) in hand, the proof of
Theorem~\ref{thm:orthog} is short.

\begin{proof}[Proof of Theorem~\ref{thm:orthog}]
Fix $M = 5$ throughout, and write $B = 1620$ for the annulus constant from
Proposition~\ref{prop:annulus}.  Let $V_j$ and $R_j$ be the local and
remainder polynomials from above.

\medskip
\noindent
\emph{Step 1: From $V_j$ to $\rho(U)$ on each circle.}\;
For $r \in [j, j{+}1]$, Proposition~\ref{prop:annulus} gives
\begin{equation}\label{eq:annulus-applied}
\int_{S^1}\lvert V_j(r\zeta) \rvert^2\,\mathrm{d}\sigma
\;\le\;
B^2\int_{S^1}\rho\bigl(V_j(r\zeta)\bigr)^2\,\mathrm{d}\sigma.
\end{equation}
Since $\rho$ is $1$-Lipschitz and $U = V_j + R_j$,
\[
\rho\bigl(V_j(z)\bigr)
\le \rho\bigl(U(z)\bigr) + \lvert R_j(z) \rvert.
\]
Squaring ($(a+b)^2 \le 2a^2 + 2b^2$) and substituting
into~\eqref{eq:annulus-applied}:
\begin{equation}\label{eq:Vj-to-U}
\int_{S^1}\lvert V_j(r\zeta) \rvert^2\,\mathrm{d}\sigma
\;\le\;
2B^2\!\int_{S^1}\rho(U(r\zeta))^2\,\mathrm{d}\sigma
\;+\; 2B^2\!\int_{S^1}\lvert R_j(r\zeta) \rvert^2\,\mathrm{d}\sigma.
\end{equation}

\medskip
\noindent
\emph{Step 2: From $U$ to $V_j$ on each circle.}\;
Applying $\lvert U \rvert^2 \le 2\lvert V_j \rvert^2 + 2\lvert R_j \rvert^2$
and~\eqref{eq:Vj-to-U}:
\begin{equation}\label{eq:circle-ineq}
\int_{S^1}\lvert U(r\zeta) \rvert^2\,\mathrm{d}\sigma
\;\le\;
4B^2\!\int_{S^1}\rho(U)^2\,\mathrm{d}\sigma
\;+\; (4B^2 + 2)\!\int_{S^1}\lvert R_j \rvert^2\,\mathrm{d}\sigma.
\end{equation}

\medskip
\noindent
\emph{Step 3: Integrate over annuli.}\;
Multiply~\eqref{eq:circle-ineq} by $2r\mathrm{e}^{-r^2}$, integrate $r$ over
$[j, j{+}1]$, and use the polar decomposition~\eqref{eq:polar}.  The
$\rho(U)$ term on the right does not depend on~$j$ (the annuli tile the
plane), while the $R_j$ term does.  Summing over $j \ge 0$:
\[
\lVert U \rVert_\mathcal{F}^2
\;\le\;
4B^2\,\lVert \rho(U) \rVert_\mathcal{F}^2
\;+\;
(4B^2 + 2)\sum_{j \ge 0}\;
\frac{1}{\pi}\int_{\{j \le \lvert z \rvert < j+1\}}\!\!
\lvert R_j \rvert^2\mathrm{e}^{-\lvert z \rvert^2}\mathrm{d} m.
\]
By Proposition~\ref{prop:total-leak} with $M = 5$, the remainder sum is at
most $\eta_5\,\lVert U \rVert_\mathcal{F}^2$.  Hence
\begin{equation}\label{eq:absorb}
\lVert U \rVert_\mathcal{F}^2
\;\le\;
4B^2\,\lVert \rho(U) \rVert_\mathcal{F}^2
\;+\;
(4B^2 + 2)\,\eta_5\,\lVert U \rVert_\mathcal{F}^2.
\end{equation}

\medskip
\noindent
\emph{Step 4: Absorption.}\;
Since $4B^2 + 2 < 10^7$ and $\eta_5 < 4 \times 10^{-11}$, the coefficient
of $\lVert U \rVert_\mathcal{F}^2$ on the right of~\eqref{eq:absorb} satisfies
$(4B^2 + 2)\,\eta_5 < 4 \times 10^{-4} < \tfrac{1}{2}$.
Moving this term to the left and using $1 - (4B^2+2)\eta_5 > 1/2$:
\[
\lVert U \rVert_\mathcal{F}^2
\;\le\; 8B^2\,\lVert \rho(U) \rVert_\mathcal{F}^2
\;=\; 8 \cdot 1620^2\,\lVert \rho(U) \rVert_\mathcal{F}^2
\;\le\; 4600^2\,\lVert \rho(U) \rVert_\mathcal{F}^2.\qedhere
\]
\end{proof}
\subsection{Generalizations of Theorem \ref{main:local}} Here, we briefly discuss some of the ideas that go into the more general Theorem \ref{thm:local_Gen}, as well as some of its associated higher dimensional generalizations. We opt to keep the discussion here brief, and refer the reader to the associated GitHub repository for detailed proofs. We also note that independent iterations of the same prompts resulted in different strategies to generalize the proof, which may be a fact of independent interest.
\medskip

\subsubsection{A special case of Theorem \ref{thm:local_Gen}} For simplicity, we discuss the special case of Theorem \ref{thm:local_Gen} where we take $k=1$ and $G=\Phi_{1,0}\coloneq \Phi_0=\overline{z}$. The general case is more technical, but follows a similar line of reasoning to what is explained below. We again refer to the associated GitHub for detailed proofs. In this special case, the result is the analogue of Theorem \ref{main:local} for the first polyanalytic fock space $\mathcal{H}_1(\mathbb{C})$. $\Phi_0$ here plays the role of the constant $1$ function in the Fock case. The proof proceeds by using the orthogonal reduction to reduce matters to proving the orthogonal coercivity bound
\begin{equation}\label{orthcoercivityhermite}
\norm{U}
\lesssim
\norm{\abs{\Phi_0 + U} - \abs{\Phi_0}}
\qquad
\bigl(U \in \mathcal{H}_1 \cap \Phi_0^\perp\bigr).
\end{equation}
This is the analogue of Theorem \ref{thm:orthog}. Once this is done, the overall architecture of the proof is the same as in the Fock space case, and with a bit of work, one can essentially reduce all of the estimates to minor variations of the ones used in this case. However, there are two minor but important points when carrying out these reductions.
\begin{enumerate}
\item The circle estimates in Section \ref{subsec:circle} only apply to trigonometric polynomials supported on positive frequencies, and also, apply to the defect function $\rho(w)=|1+w|-1$ and not $|\Phi_0+w|-\Phi_0$. Moreover, $U$ is now allowed to have a nontrivial zeroth order Fourier mode because $\Phi_0$ is supported at frequency $-1$. The fix to this simple. After writing everything in polar coordinates, one simply factors (for $r>0$) out the basis vector $\Phi_0$ on both sides of the estimate \eqref{orthcoercivityhermite}. This shifts $U$ to positive frequencies and normalizes the defect function to the one used in the Fock case, where the circle estimates in Section \ref{subsec:circle} apply.
\item The second minor difference is in the annulus localization carried out in Section \ref{subsec:blocks}. Here, instead of the monomials $z^n$ (in the Fock case), we have to understand the mass concentration properties of Hermite basis vectors $\Phi_n\coloneq \Phi_{1,n}$, to obtain the analogue of the bound \eqref{eq:monomial-bound}.
\end{enumerate}
Once these two points are addressed, the proof proceeds in an almost identical fashion to the Fock case. The proof of the general result in Theorem \ref{thm:local_Gen} is more technical, but ultimately follows a similar line of reasoning to the above (where one makes similar adjustments compared to the Fock case). We refer the reader to the GitHub for details. 

\subsubsection{Higher dimensional analogues} As mentioned earlier, Theorem \ref{main:local} has a higher dimensional analogue in $\mathbb{C}^d$. We mention here the simple analogue of this case, and refer the reader to the associated GitHub for the higher dimensional analogue of the Hermite results (which are more complicated to state). 
\medskip

To describe the setup, let $d\geq 1$ be the dimension. The analogue of the Fock norm in this case is given by
\[
\norm{F}_{\Fd}^2
:=
\frac{1}{\pi^d}\int_{\C^d} \abs{F(z)}^2 \ee^{-\abs{z}^2}\,\dd m(z),
\]
where $\Fd = \Fd(\C^d)$ is the closed subspace of entire functions with
finite norm. The analogue of Theorem \ref{main:local} is as follows.
\begin{theorem}[Local SPR at $\mathbf{1}$ in $\mathcal{F}_d(\mathbb{C}^d)$]\label{mainhigherd:local}
There exists an absolute constant $M> 0$ such
that for any $F\in \mathcal{F}_d(\mathbb{C}^d)$ we have
\[
\inf_{|\lambda|=1}\lVert F - \lambda \rVert_{\mathcal{F}_d}
\;\le\;
M\,
\bigl\|\, \lvert F \rvert - 1 \,\bigr\|_{\mathcal{F}_d}.
\]
\end{theorem}
As in the one dimensional case, this can be reduced to proving the following orthogonal coercivity estimate.
\begin{theorem}[Orthogonal coercivity on $\C^d$]\label{thm:orthogd}
Fix $d \ge 1$.  There exists a constant $C_d > 0$ such that every
$U \in \Fd$ with $U(0) = 0$ satisfies
\[
\norm{U}_{\Fd}^2
\le
C_d^2\, \norm{\rhoFn \circ U}_{\Fd}^2.
\]
\end{theorem}
This is the analogue of Theorem \ref{thm:orthog} in the higher dimensional setting. As in the one-dimensional Fock space, we still have an orthonormal basis of monomials, except here, they are of the form
\[
e_\alpha(z) := \frac{z^\alpha}{\sqrt{\alpha!}}
\]
where $\alpha\in\mathbb{N}^d$. As before, an important point in proving the orthogonal coercivity bound is in understanding the regions of space where $e_{\alpha}$ concentrates most of its mass. This allows one to carry out the analogue of the annulus frequency localization in subsection \ref{subsec:blocks}. This part is relatively similar to the one dimensional case (again, see the GitHub for the detailed proof). The part of the proof that requires a genuine additional idea is the analogue of the circle estimates in subsection \ref{subsec:circle}. This is clarified by writing $\|\cdot\|_{\mathcal{F}_d}$ norm in polar form
\begin{equation}\label{eq:polard}
\norm{F}_{\Fd}^2
=
\frac{2}{(d-1)!}
\int_0^\infty
r^{2d-1} \ee^{-r^2}
\biggl(
\int_{\Sphere} \abs{F(r\omega)}^2\,\dd \sigma_d(\omega)
\biggr)\dd r
\end{equation}
and so, one needs analogues of the circle estimates in subsection \ref{subsec:circle} (which hold for $L^2(S^1)$) for the sphere $S^{2d-1}$.

\section{Discussions on the Lean verification}\label{App}
\subsection{Lean verified statement}
This section discusses the formalization of the general statement, which can be found in the file {\scriptsize \texttt{Showcase.lean}}. This file is LLM generated and checked by the authors. At the definition level, this file depends on Mathlib only: the idea is that then, reviewing the definitions and statements in this file is enough to attest that the Lean verified statement coincides with the statement in this paper, up to the semantics of Mathlib itself. Moreover, we hope that the formulation of the statement in {\scriptsize \texttt{Showcase.lean}} is simple enough to be understood by mathematicians without deep Lean knowledge. \\
Nonetheless, we attempt to give an explanation of the content of the file. Notably, the Lean verified statement is stronger than needed: a version of the statement is proved for general elements of the space $\mathcal{H}_\kappa$, not only for finite linear combinations of elements of the base. The file {\scriptsize \texttt{Showcase.lean}} contains both the general version and the specialized version for polynomials, which is the one we will focus on.\\
We start by defining the non-normalized Hermite base: 
\begin{leancode}
    def showcaseComplexHermite (m n : Nat) (z : ℂ) : ℂ :=       Take m,n 
  Finset.sum (Finset.range (min m n + 1)) fun j =>
    ((-1 : ℂ) ^ j) * (Nat.factorial j : ℂ) *
      (Nat.choose m j : ℂ) * (Nat.choose n j : ℂ) *
      z ^ (m - j) * (star z) ^ (n - j)
\end{leancode}     
Meaning that, given $m,n\in \mathbf{N}$, we define $\Tilde{\Phi}$
\begin{equation*}
    \Tilde{\Phi}_{m,n}(z)
\end{equation*}

















\bibliographystyle{plain}
\bibliography{sources}
\end{document}
